% GENERATED FILE. Edit canonical lessons, then run npm run generate:books.
% canonical-source-hash: fnv1a64:14ae625668ca5e60
% canonical-lessons: IT-C528-serie, IT-C528-puntata, IT-C528-episodio, IT-C528-spoiler, IT-C528-telecomando

\chapter{Una serie da vedere}
\label{ch:it-una-serie-da-vedere}
\hlchaptermodality{528}

\begin{chapteropening}
\textbf{By the end of this chapter:} \emph{I can talk about watching TV series.}

Complete the last lesson of chapter 528: I can talk about watching TV series.
\end{chapteropening}

\section[la serie]{la serie — series}

\label{lesson:IT-C528-serie}



\begin{quote}
Before the new one: say the Italian for comic, then the Italian for poet.
\end{quote}

\subsection*{You'll want to know: la serie}
\textbf{la serie} — \textquotedblleft{}series\textquotedblright{}.

\begin{quote}
Stiamo guardando una serie poliziesca.
\end{quote}

\begin{quote}
\emph{We are watching a detective series.}
\end{quote}

The same form in the plural: \emph{le serie}. \emph{Una serie TV}.

\begin{cousinweb}
\emph{serie} is a learned borrowing from Latin \emph{seriēs} \textquotedblleft{}a row, a chain\textquotedblright{} (Lewis \& Short), from \emph{serere} \textquotedblleft{}to join\textquotedblright{}.
\end{cousinweb}

\begin{grammarlens}[title={What you've built}]
One more word for telling, explaining and describing at B1.
\end{grammarlens}

\subsection*{Your turn}
\begin{itemize}
\raggedright
  \item \emph{Say it:} \emph{la serie}
  \item \emph{Say it:} \emph{la serie}, once more
  \item \emph{Say it:} a sentence of your own with \emph{la serie}
  \item \emph{Recall:} say the Italian for comic, then the Italian for poet, then say \emph{la serie} again
\end{itemize}

\begin{tcolorbox}[breakable,colback=teal!4,colframe=teal!35!black,title={Before you move on}]
What does la serie mean? (\textquotedblleft{}Series\textquotedblright{}.) Say it once more.
\end{tcolorbox}
\section[la puntata]{la puntata — instalment, episode (of a series)}

\label{lesson:IT-C528-puntata}



\begin{quote}
Before the new one: say the Italian for poet, then the Italian for series.
\end{quote}

\subsection*{You'll want to know: la puntata}
\textbf{la puntata} — \textquotedblleft{}instalment, episode (of a series)\textquotedblright{}.

\begin{quote}
Stasera c'è l'ultima puntata.
\end{quote}

\begin{quote}
\emph{Tonight is the last episode.}
\end{quote}

Built on \emph{punto}; \emph{a puntate} — in instalments.

\begin{cousinweb}
\emph{puntata} is built on \emph{punto} \textquotedblleft{}point\textquotedblright{}, inherited from Latin \emph{punctum} (Lewis \& Short), from \emph{pungere} \textquotedblleft{}to prick\textquotedblright{}.
\end{cousinweb}

\begin{grammarlens}[title={What you've built}]
One more word for telling, explaining and describing at B1.
\end{grammarlens}

\subsection*{Your turn}
\begin{itemize}
\raggedright
  \item \emph{Say it:} \emph{la puntata}
  \item \emph{Say it:} \emph{la puntata}, once more
  \item \emph{Say it:} a sentence of your own with \emph{la puntata}
  \item \emph{Recall:} say the Italian for poet, then the Italian for series, then say \emph{la puntata} again
\end{itemize}

\begin{tcolorbox}[breakable,colback=teal!4,colframe=teal!35!black,title={Before you move on}]
What does la puntata mean? (\textquotedblleft{}Instalment, episode (of a series)\textquotedblright{}.) Say it once more.
\end{tcolorbox}
\section[l'episodio]{l'episodio — episode; incident}

\label{lesson:IT-C528-episodio}



\begin{quote}
Before the new one: say the Italian for series, then the Italian for instalment.
\end{quote}

\subsection*{You'll want to know: l'episodio}
\textbf{l'episodio} — \textquotedblleft{}episode; incident\textquotedblright{}.

\begin{quote}
Ho visto tre episodi di fila.
\end{quote}

\begin{quote}
\emph{I watched three episodes in a row.}
\end{quote}

Also an incident: \emph{un brutto episodio}.

\begin{cousinweb}
\emph{episodio} is a learned borrowing from Greek \emph{epeisódion} \textquotedblleft{}a part added between choral songs\textquotedblright{} (EtymDB, under English \emph{episode}), from \emph{eísodos} \textquotedblleft{}an entrance\textquotedblright{}.
\end{cousinweb}

\begin{grammarlens}[title={What you've built}]
One more word for telling, explaining and describing at B1.
\end{grammarlens}

\subsection*{Your turn}
\begin{itemize}
\raggedright
  \item \emph{Say it:} \emph{l'episodio}
  \item \emph{Say it:} \emph{l'episodio}, once more
  \item \emph{Say it:} a sentence of your own with \emph{l'episodio}
  \item \emph{Recall:} say the Italian for series, then the Italian for instalment, then say \emph{l'episodio} again
\end{itemize}

\begin{tcolorbox}[breakable,colback=teal!4,colframe=teal!35!black,title={Before you move on}]
What does l'episodio mean? (\textquotedblleft{}Episode; incident\textquotedblright{}.) Say it once more.
\end{tcolorbox}
\section[lo spoiler]{lo spoiler — spoiler}

\label{lesson:IT-C528-spoiler}



\begin{quote}
Before the new one: say the Italian for instalment, then the Italian for episode.
\end{quote}

\subsection*{You'll want to know: lo spoiler}
\textbf{lo spoiler} — \textquotedblleft{}spoiler\textquotedblright{}.

\begin{quote}
Non dirmi niente, niente spoiler!
\end{quote}

\begin{quote}
\emph{Do not tell me anything, no spoilers!}
\end{quote}

From English; the plural is \emph{gli spoiler}.

\begin{cousinweb}
\emph{spoiler} was borrowed from English \emph{spoiler}, built on \emph{spoil}, from Latin \emph{spoliāre} \textquotedblleft{}to strip\textquotedblright{} (Lewis \& Short).
\end{cousinweb}

\begin{grammarlens}[title={What you've built}]
One more word for telling, explaining and describing at B1.
\end{grammarlens}

\subsection*{Your turn}
\begin{itemize}
\raggedright
  \item \emph{Say it:} \emph{lo spoiler}
  \item \emph{Say it:} \emph{lo spoiler}, once more
  \item \emph{Say it:} a sentence of your own with \emph{lo spoiler}
  \item \emph{Recall:} say the Italian for instalment, then the Italian for episode, then say \emph{lo spoiler} again
\end{itemize}

\begin{tcolorbox}[breakable,colback=teal!4,colframe=teal!35!black,title={Before you move on}]
What does lo spoiler mean? (\textquotedblleft{}Spoiler\textquotedblright{}.) Say it once more.
\end{tcolorbox}
\section[il telecomando]{il telecomando — remote control}

\label{lesson:IT-C528-telecomando}



\begin{quote}
Before the new one: say the Italian for episode, then the Italian for spoiler.
\end{quote}

\subsection*{You'll want to know: il telecomando}
\textbf{il telecomando} — \textquotedblleft{}remote control\textquotedblright{}.

\begin{quote}
Dov'è il telecomando? Non lo trovo.
\end{quote}

\begin{quote}
\emph{Where is the remote? I cannot find it.}
\end{quote}

\emph{Tele-} \textquotedblleft{}far\textquotedblright{} as in \emph{televisione}.

\begin{cousinweb}
\emph{telecomando} is built from Greek \emph{tēle} \textquotedblleft{}far\textquotedblright{} and \emph{comando} \textquotedblleft{}command\textquotedblright{}, from Latin \emph{mandāre} \textquotedblleft{}to entrust\textquotedblright{} (Lewis \& Short).
\end{cousinweb}

\begin{grammarlens}[title={What you've built}]
One more word for telling, explaining and describing at B1.
\end{grammarlens}

\subsection*{Your turn}
\begin{itemize}
\raggedright
  \item \emph{Say it:} \emph{il telecomando}
  \item \emph{Say it:} \emph{il telecomando}, once more
  \item \emph{Say it:} a sentence of your own with \emph{il telecomando}
  \item \emph{Recall:} say the Italian for episode, then the Italian for spoiler, then say \emph{il telecomando} again
\end{itemize}

\begin{tcolorbox}[breakable,colback=teal!4,colframe=teal!35!black,title={Before you move on}]
What does il telecomando mean? (\textquotedblleft{}Remote control\textquotedblright{}.) Say it once more.
\end{tcolorbox}
